% Options for packages loaded elsewhere
\PassOptionsToPackage{unicode}{hyperref}
\PassOptionsToPackage{hyphens}{url}
\documentclass[
  spanish,
  12pt,
  letterpaper,
]{article}
\usepackage{xcolor}
\usepackage[margin=2.4cm]{geometry}
\usepackage{amsmath,amssymb}
\setcounter{secnumdepth}{-\maxdimen} % remove section numbering
\usepackage{iftex}
\ifPDFTeX
  \usepackage[T1]{fontenc}
  \usepackage[utf8]{inputenc}
  \usepackage{textcomp} % provide euro and other symbols
\else % if luatex or xetex
  \usepackage{unicode-math} % this also loads fontspec
  \defaultfontfeatures{Scale=MatchLowercase}
  \defaultfontfeatures[\rmfamily]{Ligatures=TeX,Scale=1}
\fi
\usepackage{lmodern}
\ifPDFTeX\else
  % xetex/luatex font selection
  \setmainfont[]{Times New Roman}
\fi
% Use upquote if available, for straight quotes in verbatim environments
\IfFileExists{upquote.sty}{\usepackage{upquote}}{}
\IfFileExists{microtype.sty}{% use microtype if available
  \usepackage[]{microtype}
  \UseMicrotypeSet[protrusion]{basicmath} % disable protrusion for tt fonts
}{}
\usepackage{setspace}
\makeatletter
\@ifundefined{KOMAClassName}{% if non-KOMA class
  \IfFileExists{parskip.sty}{%
    \usepackage{parskip}
  }{% else
    \setlength{\parindent}{0pt}
    \setlength{\parskip}{6pt plus 2pt minus 1pt}}
}{% if KOMA class
  \KOMAoptions{parskip=half}}
\makeatother
\ifLuaTeX
\usepackage[bidi=basic,shorthands=off,]{babel}
\else
\usepackage[bidi=default,shorthands=off,]{babel}
\fi
\ifPDFTeX
\else
\babelfont{rm}[]{Times New Roman}
\fi
\ifLuaTeX
  \usepackage{selnolig} % disable illegal ligatures
\fi
\setlength{\emergencystretch}{3em} % prevent overfull lines
\providecommand{\tightlist}{%
  \setlength{\itemsep}{0pt}\setlength{\parskip}{0pt}}

\usepackage{titlesec}
\titleformat{\section}{\centering\bfseries}{\thesection}{1em}{}
\titleformat{\subsection}{\raggedright\bfseries}{\thesubsection}{1em}{}
\titleformat{\subsubsection}{\raggedright\bfseries\itshape}{\thesubsubsection}{1em}{}
\titleformat{\paragraph}[runin]{\raggedright\bfseries}{\theparagraph}{1em}{}[.]
\titleformat{\subparagraph}[runin]{\raggedright\bfseries\itshape}{\thesubparagraph}{1em}{}
\titlespacing*{\paragraph}{0cm}{*1}{1em}
\titlespacing*{\subparagraph}{0cm}{*1}{1em}

\setlength{\parindent}{0cm}
\setlength{\parskip}{0pt}

\usepackage{caption}
\captionsetup[table]{labelfont=bf,textfont=it,labelsep=newline,justification=raggedright,singlelinecheck=false}
\captionsetup[figure]{labelfont=bf,textfont=it,labelsep=newline,justification=raggedright,singlelinecheck=false}

\usepackage{xurl}
\tolerance=9999
\emergencystretch=3em
\hbadness=10000

\usepackage{bookmark}
\IfFileExists{xurl.sty}{\usepackage{xurl}}{} % add URL line breaks if available
\urlstyle{same}
\hypersetup{
  pdflang={es},
  hidelinks,
  pdfcreator={LaTeX via pandoc}}

\author{}
\date{}

\begin{document}

\setstretch{1}
\begin{center}

\textbf{RECOLECCION DE DATOS MEDIANTE ENCUESTAS}

\emph{D. A. Flores Cordero}

7690-14-745 Universidad Mariano Gálvez

047 Seminario de Tecnologías de Información

\emph{dfloresc3@miumg.edu.gt}

\end{center}

\section{Resumen}\label{resumen}

Este trabajo investiga los criterios necesarios para diseñar una
encuesta efectiva y óptima, con el objetivo de aplicarlos en la
construcción de un instrumento propio en una herramienta en línea. La
investigación partió de tres fuentes especializadas en metodología de
encuestas y experiencia de usuario, contrastadas con un análisis previo
sobre los elementos que determinan la calidad de un cuestionario: la
definición del objetivo, la delimitación del público, el tipo de
pregunta según el propósito de la encuesta y la lógica condicional entre
preguntas. A partir de ese análisis se identificaron cinco criterios
centrales ---objetivo claro, brevedad, lenguaje neutro y comprensible,
orden lógico de las preguntas y transparencia sobre el uso de la
información--- que se aplicaron directamente al diseñar una encuesta de
prueba en una herramienta en línea. Los resultados confirman que la
calidad de una encuesta depende menos de la plataforma utilizada y más
de las decisiones tomadas antes de escribir la primera pregunta: sin un
objetivo definido, ningún criterio de redacción compensa la falta de
dirección del instrumento. Se concluye que definir el objetivo y el
público desde el inicio, junto con la brevedad y la neutralidad del
lenguaje, son los factores que más influyen en la tasa de respuesta y en
la calidad de los datos recolectados.

\textbf{Palabras claves:} encuestas, diseño de cuestionarios,
recolección de datos, experiencia de usuario

\section{Desarrollo del tema}\label{desarrollo-del-tema}

Una encuesta es un método para obtener información directamente de los
sujetos de estudio sobre sus opiniones, actitudes o comportamientos
(Galán Amador, s.f.). A diferencia de otros instrumentos de
investigación, su valor no depende de la herramienta con la que se
construye, sino de las decisiones que se toman antes de redactar la
primera pregunta.

\subsection{Definir el objetivo antes de
diseñar}\label{definir-el-objetivo-antes-de-diseuxf1ar}

Todo diseño de encuesta debe partir de una pregunta previa: qué se
quiere descubrir y qué límites tiene esa búsqueda. QuestionPro (s.f.)
coincide en que el objetivo claro es el criterio más importante para
construir un buen cuestionario, porque de él dependen el tipo de
muestreo, la cantidad de preguntas y el uso que se le dará a la
información recolectada. Sin este paso, ningún criterio de redacción
posterior compensa la falta de dirección del instrumento.

\subsection{Delimitar el público y ser
transparente}\label{delimitar-el-puxfablico-y-ser-transparente}

Definido el objetivo, corresponde elegir a quién se dirige la encuesta e
informarle por qué se le pregunta. Lucid (s.f.) recomienda ser
transparente sobre cómo se usarán los datos recolectados y mantener las
respuestas anónimas cuando corresponda, ya que esa claridad incrementa
la disposición del encuestado a responder con honestidad.

\subsection{El tipo de pregunta según el
propósito}\label{el-tipo-de-pregunta-seguxfan-el-propuxf3sito}

El tipo de pregunta no es una elección arbitraria: depende del objetivo
definido al inicio. Si la encuesta es exploratoria, conviene usar
preguntas abiertas, que permiten conocer cómo percibe el encuestado un
problema o una solución; si la encuesta busca datos comparables y
fáciles de procesar, corresponden preguntas cerradas o de opción
múltiple (Lucid, s.f.). La cantidad de preguntas también varía según el
objetivo --- no existe un número ideal fijo, sino el mínimo necesario
para cubrirlo.

\subsection{Preguntas breves, claras y
neutrales}\label{preguntas-breves-claras-y-neutrales}

Tanto QuestionPro (s.f.) como Lucid (s.f.) coinciden en tres criterios
de redacción: preguntas breves (no más de 20 palabras), lenguaje simple
y comprensible, y neutralidad para evitar sesgar la respuesta o
incomodar al encuestado. Preguntar ``¿cuánto calcula que gastó?'' en
lugar de ``¿cuánto gastó?'' es un ejemplo de cómo una pregunta mal
calibrada obtiene respuestas menos confiables que una pregunta que
reconoce los límites de la memoria del encuestado.

\subsection{Orden lógico y lógica
condicional}\label{orden-luxf3gico-y-luxf3gica-condicional}

No todas las preguntas aplican a todos los encuestados: la respuesta a
una pregunta puede abrir o cerrar el paso a otras. Además, las preguntas
deben agruparse por tema y seguir un orden lógico --- pasar de un tema a
otro sin transición confunde al encuestado y afecta la calidad de la
respuesta (Lucid, s.f.).

\subsection{Aplicación práctica}\label{aplicaciuxf3n-pruxe1ctica}

Con base en estos criterios se diseñó una encuesta de prueba en una
herramienta en línea, disponible en:
\url{https://form.jotform.com/262052207277049}. El instrumento aplica el
objetivo claro, la brevedad y el orden lógico descritos en los puntos
anteriores como criterios de diseño.

\section{Observaciones y comentarios}\label{observaciones-y-comentarios}

El ejercicio confirmó que ninguna herramienta en línea compensa un
objetivo mal definido: la tecnología solo ejecuta lo que el investigador
ya decidió de antemano. También quedó claro que la brevedad y la
claridad del lenguaje pesan más que la cantidad de preguntas --- una
encuesta larga y bien intencionada puede fallar simplemente por cansar
al encuestado. El mayor reto no fue elegir la herramienta, sino resistir
la tentación de preguntar de más.

\section{Conclusiones}\label{conclusiones}

\begin{enumerate}
\def\labelenumi{\arabic{enumi}.}
\tightlist
\item
  El objetivo de la encuesta determina todas las decisiones posteriores:
  público, tipo de pregunta y extensión.
\item
  La brevedad y el lenguaje neutro aumentan la probabilidad de obtener
  respuestas completas y confiables.
\item
  La lógica condicional entre preguntas evita exponer al encuestado a
  preguntas que no le corresponden.
\item
  La calidad de una encuesta depende del diseño previo, no de la
  plataforma utilizada para construirla.
\end{enumerate}

\section{Bibliografía}\label{bibliografuxeda}

Galán Amador, M. (s.f.). \emph{La encuesta} {[}entrada de blog{]}. Guía
metodológica para investigación.
http://manuelgalan.blogspot.com/p/guia-metodologica-para-investigacion.html

Lucid. (s.f.). \emph{Consejos para crear encuestas de usuarios
efectivas}.
https://lucid.co/es/blog/consejos-para-crear-encuestas-de-usuarios-efectivas

QuestionPro. (s.f.). \emph{Criterios para elaborar encuestas de forma
efectiva}.
https://www.questionpro.com/blog/es/criterios-para-elaborar-encuestas/

\begin{center}\rule{0.5\linewidth}{0.5pt}\end{center}

Repositorio del documento (código fuente LaTeX):
\url{https://github.com/dan-desa/seminario-tecnologias-2026/tree/main/foro-01}

\end{document}
